% Standalone problem document, generated from the adjacent README.md.
% Compile with XeLaTeX. Mathematical content is maintained in Markdown.
\documentclass[11pt,a4paper]{article}
\usepackage[fontsize=10.5pt]{scrextend}
\usepackage[margin=22mm,headheight=15pt,headsep=9mm,footskip=11mm]{geometry}
\usepackage{fontspec}
\setmainfont{texgyrepagella}[Extension=.otf,UprightFont=*-regular,BoldFont=*-bold,ItalicFont=*-italic,BoldItalicFont=*-bolditalic]
\setsansfont{texgyreheros}[Extension=.otf,UprightFont=*-regular,BoldFont=*-bold,ItalicFont=*-italic,BoldItalicFont=*-bolditalic]
\setmonofont{texgyrecursor}[Extension=.otf,UprightFont=*-regular,BoldFont=*-bold,ItalicFont=*-italic,BoldItalicFont=*-bolditalic,Scale=MatchLowercase]
\usepackage{amsmath,amssymb}
\usepackage{unicode-math}
\setmathfont{texgyrepagella-math.otf}
\usepackage{xcolor,microtype,enumitem,needspace,fancyhdr}
\usepackage{bookmark}
\definecolor{ink}{HTML}{17344B}
\definecolor{accent}{HTML}{197D80}
\definecolor{muted}{HTML}{596773}
\hypersetup{colorlinks=true,linkcolor=accent,urlcolor=accent,pdftitle={MI-07 - The
triangle conjecture for the maximal symmetric
modulus},pdfauthor={Open Problems in Numerical Linear Algebra}}
\urlstyle{same}
\setlength{\parindent}{0pt}
\setlength{\parskip}{5pt plus 1pt minus 1pt}
\linespread{1.04}
\setlength{\emergencystretch}{3em}
\widowpenalty=10000
\clubpenalty=10000
\displaywidowpenalty=10000
\allowdisplaybreaks[1]
\setlist{nosep,leftmargin=1.5em}
\providecommand{\tightlist}{\setlength{\itemsep}{0pt}\setlength{\parskip}{0pt}}
\providecommand{\pandocbounded}[1]{#1}
\pagestyle{fancy}
\fancyhf{}
\fancyhead[L]{\sffamily\footnotesize\color{muted}OPEN PROBLEMS IN NUMERICAL LINEAR ALGEBRA}
\fancyhead[R]{\sffamily\bfseries\footnotesize\color{accent}MI-07}
\fancyfoot[L]{\parbox[b]{0.9\textwidth}{\sffamily\fontsize{8}{10}\selectfont\color{muted}Editorial ratings; a bounded search cannot certify that a problem remains unsolved.\\Literature check: 2026-09-12}}
\fancyfoot[R]{\sffamily\footnotesize\color{muted}\thepage}
\renewcommand{\headrulewidth}{0.3pt}
\renewcommand{\headrule}{\hbox to\headwidth{\color{accent}\leaders\hrule height \headrulewidth\hfill}}
\makeatletter
\renewcommand{\section}{\Needspace{5\baselineskip}\@startsection{section}{1}{0pt}{12pt plus 2pt minus 2pt}{4pt}{\sffamily\bfseries\large\color{ink}}}
\renewcommand{\subsection}{\Needspace{4\baselineskip}\@startsection{subsection}{2}{0pt}{9pt plus 1pt minus 1pt}{3pt}{\sffamily\bfseries\normalsize\color{ink}}}
\makeatother
\setcounter{secnumdepth}{0}
\begin{document}
{\sffamily\small\color{muted}Matrix inequalities and norms\par}
\vspace{3pt}
{\raggedright\hyphenpenalty=10000\exhyphenpenalty=10000\sffamily\bfseries\fontsize{21}{24}\selectfont\color{ink}The
triangle conjecture for the maximal symmetric modulus\par}
\vspace{7pt}
{\sffamily\small\textbf{Difficulty:} challenging\quad\textbf{Importance:} interesting
to specialist\par}
{\sffamily\small\bfseries\color{accent}Status: Lean verified\par}
\vspace{4pt}
{\color{accent}\hrule height 0.8pt}
\vspace{6pt}
\textbf{Rating rationale:} Combining spectral-order suprema with
ordinary matrix-order domination requires new analysis; the immediate
audience is operator-inequality specialists.

\subsection{Lean proof and verification evidence ---
2026-09-12}\label{lean-proof-and-verification-evidence-2026-09-12}

\textbf{The original constant-one MI-07 conjecture is false, with a
complete Lean-verified counterexample.} At the fixed rational pair

\[A=\begin{pmatrix}1&0\\0&0\end{pmatrix},\qquad
B=\begin{pmatrix}0&5/12\\0&0\end{pmatrix},\]

the actual maximal moduli are \(M(A)=A\), \(M(B)=(5/12)I\) and
\(M(A+B)=(13/12)I\). The left trace is \(13/6\), whereas every complex
unitary pair gives right trace \(11/6\). The right-minus-left trace is
\(-1/3\), which excludes ordinary PSD domination. The
\href{https://github.com/sgstepaniants/OpenProblemsInNLA/tree/f55777156432043de4201747a3759e0c6485e568/matrix-inequalities-and-norms/MI-07/lean}{proof
at revision f557771} retains the full original complex-matrix and
unitary quantifiers and proves the actual root-sequence limits at all
three counterexample arguments.

\textbf{Mathematical counterexample and informal proof:} Matthew J.
Colbrook, Department of Applied Mathematics and Theoretical Physics,
University of Cambridge. \textbf{Lean formalization:} George
Stepaniants, Department of Computing and Mathematical Sciences,
California Institute of Technology, Pasadena, California, USA, with
AI-agent assistance.

The seven checked declarations in
\href{https://github.com/sgstepaniants/OpenProblemsInNLA/blob/f55777156432043de4201747a3759e0c6485e568/matrix-inequalities-and-norms/MI-07/lean/Solution.lean}{Solution.lean}
are:

\begin{itemize}
\tightlist
\item
  \textit{NLA.MI07.modulus\_eq\_sqrt}: the genuine positive square-root
  modulus.
\item
  \textit{NLA.MI07.maximalModulus\_eq\_of\_tendsto}: identification of
  the actual limit from proved convergence.
\item
  \textit{NLA.MI07.root\_limit\_iff\_spectralNorm}: equivalence with
  convergence in the genuine Euclidean operator norm.
\item
  \textit{NLA.MI07.witness\_moduli}: all exact polar factors and
  positivity facts.
\item
  \textit{NLA.MI07.witness\_root\_limits}: the three actual
  root-sequence limits.
\item
  \textit{NLA.MI07.counterexample}: the exact all-unitary trace
  obstruction.
\item
  \textit{NLA.MI07.not\_triangleConjecture}: negation of the complete
  original universal assertion.
\end{itemize}

Two independent agents reviewed the
\href{https://github.com/ajt60gaibb/OpenProblemsInNLA/blob/main/matrix-inequalities-and-norms/MI-07/lean/reviews}{frozen
statements and completed proof}.
\href{https://github.com/sgstepaniants/OpenProblemsInNLA/actions/runs/34709291624}{Linux
run 34709291624} matched all seven exports with the sandboxed Comparator
and replayed them through Lean's default kernel. The
\href{https://github.com/ajt60gaibb/OpenProblemsInNLA/blob/main/matrix-inequalities-and-norms/MI-07/lean/verification/linux-2026-09-12}{original
artifacts and operational audit} retain source hashes, isolation checks
and rejection controls. The
\href{https://github.com/ajt60gaibb/OpenProblemsInNLA/blob/main/matrix-inequalities-and-norms/MI-07/lean/reviews/proof-axioms.log}{transitive
axiom reports} contain only \textit{propext}, \textit{Classical.choice}
and \textit{Quot.sound}. These are independent agent reviews, not
external human peer review.

The proof pins \textbf{Lean 4.33.1},
\href{https://github.com/alerad/leancert/tree/621a43d7cf21f87872392a01e874f2f1dbddc926}{LeanCert
621a43d} and
\href{https://github.com/leanprover-community/mathlib4/tree/0df444a360eaa60ab8c11dca51a86af692955474}{Mathlib
0df444a}. See the
\href{https://github.com/ajt60gaibb/OpenProblemsInNLA/blob/main/matrix-inequalities-and-norms/MI-07/lean/formalization.yaml}{manifest},
\href{https://github.com/ajt60gaibb/OpenProblemsInNLA/blob/main/matrix-inequalities-and-norms/MI-07/lean/lake-manifest.json}{dependency
pins} and
\href{https://github.com/ajt60gaibb/OpenProblemsInNLA/blob/main/matrix-inequalities-and-norms/MI-07/lean/NUMERICAL_TARGETS.md}{numerical
targets}. On a documented
\href{https://github.com/ajt60gaibb/OpenProblemsInNLA/blob/main/tools/lean/HARNESS.md}{non-root
Linux host}, reproduce from the verified revision with:

\begin{verbatim}
tools/lean/bootstrap.sh /absolute/path/to/nla-lean-tools
tools/lean/selftest.sh /absolute/path/to/nla-lean-tools
tools/lean/verify.sh \
  matrix-inequalities-and-norms/MI-07/lean \
  /absolute/path/to/nla-lean-tools
\end{verbatim}

The manuscript's stronger no-finite-constant theorem and generic
convergence for unrelated matrices are outside these seven exports. The
original informal proof and its historical audit remain below.

\subsection{Resolution --- 2026-09-11}\label{resolution-2026-09-11}

\textbf{Negative result by Matthew J. Colbrook}, Department of Applied
Mathematics and Theoretical Physics, University of Cambridge.
\textbf{Independent proof review: PASS.}

No finite two-unitary domination constant exists for the maximal
symmetric modulus, already in dimension two. The rank-one family gives
the necessary bound \(C\ge\sqrt{1+t^2}/t\) for every \(t>0\).

The exact target is resolved. The original statement and source evidence
are retained below; its former difficulty rating is historical.

\textbf{Primary manuscript:}
\href{https://github.com/ajt60gaibb/OpenProblemsInNLA/blob/main/matrix-inequalities-and-norms/MI-07/solution.pdf}{complete
proof PDF},
\href{https://github.com/ajt60gaibb/OpenProblemsInNLA/blob/main/matrix-inequalities-and-norms/MI-07/solution.tex}{standalone
TeX}, Theorem 1.1 and its proof;
\href{https://github.com/ajt60gaibb/OpenProblemsInNLA/blob/main/matrix-inequalities-and-norms/MI-07/solution.md}{authorship
and scope}. The
\href{https://github.com/ajt60gaibb/OpenProblemsInNLA/blob/main/references/colbrook-matrix-2026-09-11/verification/reviews/MI-07-review.md}{independent
review} checks the full original argument and records its hash. The
draft was AI-assisted; the 2026-09-11 review was independent agent
verification, without external human peer review or formal
certification. The later Lean verification above covers the complete
original constant-one target.
\href{https://github.com/ajt60gaibb/OpenProblemsInNLA/blob/main/references/colbrook-matrix-2026-09-11/README.md}{Submission
record}.

\subsection{Problem statement}\label{problem-statement}

For a complex square matrix \(X\), let \(|X|=(X^*X)^{1/2}\) and define
its maximal symmetric modulus by the finite-dimensional limit

\[M(X)=\lim_{r\to\infty}\bigl(|X|^r+|X^*|^r\bigr)^{1/r},\qquad r\in\mathbb N.\]

This is the supremum of \(|X|\) and \(|X^*|\) in Olson's spectral order.
For every \(n\ge1\) and every \(A,B\in\mathbb C^{n\times n}\), must
there exist unitary \(U,V\) of order \(n\) satisfying

\[M(A+B)\preceq U M(A)U^*+V M(B)V^*?\]

The order in this displayed inequality is ordinary positive-semidefinite
order.

\subsection{Why it matters}\label{why-it-matters}

The modulus combines both positive polar factors through spectral order.
A sharp triangle inequality would give a new way to control positive
representatives of nonnormal matrix sums.

\newpage
\subsection{References}\label{references}

\begin{enumerate}
\def\labelenumi{\arabic{enumi}.}
\tightlist
\item
  J.-C. Bourin and E.-Y. Lee, \emph{Averages over matrix unitary orbits
  and spectral order}, arXiv:2606.15624v2 (18 June 2026), §§3.1--3.3,
  Question 3.11 and Conjecture 3.12.
  \href{https://arxiv.org/html/2606.15624}{Primary text}.
\item
  J.-C. Bourin and E.-Y. Lee, \emph{Some hybrid matrix triangle
  inequalities}, arXiv:2606.29188v1 (28 June 2026), §3, Lemma 3.3 and
  Theorem 3.4. \href{https://arxiv.org/html/2606.29188}{Primary text}.
\end{enumerate}

\subsection{Status check --- 2026-09-10}\label{status-check-2026-09-10}

Latest versions v2 and v1 were checked. The follow-up proves
inequalities with the ordinary modulus of the total sum on the left,
while the present conjecture has its maximal symmetric modulus there.
Its conclusion therefore does not establish the displayed assertion.
Searches included \textit{maximal symmetric modulus conjecture},
\textit{Bourin Lee 2606.15624 conjecture}, and
\textit{Some hybrid matrix triangle inequalities}. No later
resolution was located.

\textbf{Audit update (2026-09-10):} Rechecked Bourin--Lee Conjecture
3.12 and their hybrid follow-up, then searched for maximal-modulus
triangle results. The follow-up does not replace its left-hand ordinary
modulus by the maximal modulus required here. This is a bounded
literature check, not a proof that no solution exists.
\end{document}
